% -----------------------------------------------------------------------------
% This file is automatically generated.
% Do not edit manually.
% Source: notebooks/support/probability-distributions/support.Rmd
% -----------------------------------------------------------------------------



\noindent We see that the result for the cumulative function equals that of the sum of the two individual probabilities.

\noindent Probabilities in the right-hand tail $P(k > k)$ can be calculated by passing the additional argument \ttblue{lower.tail = FALSE}. For example, the probability $P(k > 1)$ is calculated by


\par\addvspace{\topsep}
\begingroup
\fvset{listparameters={%
  \setlength{\topsep}{0pt}%
  \setlength{\partopsep}{0pt}%
  \setlength{\parsep}{0pt}%
  \setlength{\itemsep}{0pt}%
}}
\begin{Verbatim}[breaklines=true,formatcom=\color{ada_blue}]
hypergeom.sf(1, M=(17) + (314), n=17, N=60)
\end{Verbatim}
\begin{Verbatim}[breaklines=true,formatcom=\color{ada_light_blue}]
[1] 0.8481751
\end{Verbatim}
\endgroup
\par\addvspace{\topsep}
